\section{Surge attribution: background, shock and contagion}\label{sec:attribution}

A salient external event is often followed by a rise in recorded incidents. Three
questions then arise. How many incidents in the following days would have occurred
without the event? How many respond directly to it? How many were triggered by other
incidents in the cascade that the event set off? Under model \eqref{eq:th-intensity}
these questions have exact answers in terms of the parameters, because the cluster
representation (\cref{sec:th-cluster}) assigns every incident to exactly one
immigrant ancestor.

\Cref{thm:attribution} splits the expected count in a post-event window into
background, direct and contagion components, \cref{prop:irf} gives closed forms for
exponential kernels, and \cref{prop:at-channel} isolates the part of the contagion that
passes from online to offline streams. All of these quantities are functionals of the
model parameters. They describe what the fitted model implies, and they inherit the
identification limits of \cref{sec:identification}: a component labelled ``contagion'' is
contagion only if the model is correct, including its background and its list of events
(\cref{rem:at-causal}). The ingredients are classical: the resolvent and the first-moment
formula go back to \citet{hawkes1971} (see also \citealp{bacry2015,laub2021}), the split by
immigrant ancestor is the cluster representation of \citet{hawkes1974}, and resolvents of
Volterra equations are treated in \citet{gripenberg1990}. We have not found the
hypoexponential closed forms of \cref{prop:irf}(d) or the channel decomposition stated in
this form, although both are short consequences of classical identities.
\Cref{app:attrproofs} gives the proofs.

\subsection{Set-up}\label{sec:at-setup}

Throughout this section \cref{ass:th-model} holds and $\spr(\Amat)<1$. The process
starts from the empty history at time $0$, as in \eqref{eq:th-intensity}. Fix an
event $e$ with time $\tau_e\ge0$ and size $S_e$, a horizon $H\in(0,\infty]$ and the
post-event window $W_e=(\tau_e,\tau_e+H]$. Inequalities between vectors or matrices
hold entrywise. We collect the loadings into the \emph{direct-response vector}
\[
\bm d_e=S_e\,(\upsilon_1,\dots,\upsilon_M)^{\top}.
\]
Its $k$th entry is the expected number of stream-$k$ incidents that respond directly to
event $e$. The immigrant intensities of \cref{sec:th-cluster} split by source into
\[
\bm\mu(t),\qquad
\bm\nu_e(t)=S_e\bigl(\upsilon_k\psi_k(t-\tau_e)\bigr)_{k\in[M]},\qquad
\bm\nu_{-e}(t)=\sum_{e'\neq e}\bm\nu_{e'}(t).
\]
For a matrix function $\bm F$ and a vector function $\bm f$ on $[0,\infty)$, let
$(\bm F*\bm f)(t)=\int_0^t\bm F(t-s)\bm f(s)\,\dd s$. Matrix convolution is defined in
the same way, with $\Phik^{*1}=\Phik$ and $\Phik^{*(m+1)}=\Phik*\Phik^{*m}$. Write
$\bar\Rker(u)=\int_0^u\Rker$. For $\tau\ge0$, the \emph{cascade operator} maps an
immigrant intensity to the mean intensity it generates:
\begin{equation}\label{eq:at-cascade}
  (\mathcal K_\tau\bm f)(t)=\bm f(t)+\int_\tau^t\Rker(t-s)\,\bm f(s)\,\dd s,
  \qquad t>\tau,\qquad \mathcal K:=\mathcal K_0 .
\end{equation}

Since $\spr(\Amat)<1$, the series $\Rker=\sum_{m\ge1}\Phik^{*m}$ converges in
$L^1[0,\infty)$ with $\int_0^\infty\Rker=\Mmat-\Imat$. In the cluster representation the
incidents of generation $m$ have mean intensity $\Phik^{*m}*\bm\nu$, so an immigrant
intensity $\bm\nu$ generates the mean intensity $\mathcal K\bm\nu$, the unique locally
integrable solution of the renewal equation $\bm m=\bm\nu+\Phik*\bm m$
(\cref{lem:at-mean} in \cref{app:at-mean}).

\subsection{The decomposition}\label{sec:at-decomp}

\begin{theorem}[Surge attribution]\label{thm:attribution}
Under the standing assumptions of \cref{sec:at-setup}, let
$\Psi_k(H)=\int_0^H\psi_k$.
\begin{enumerate}[(a)]
  \item \emph{Decomposition.} For every stream $k$,
  \begin{equation}\label{eq:at-decomp}
    \E\,N_k(W_e)=B_k+D_k+C_k ,
  \end{equation}
  where
  \begin{align}
    D_k&=\upsilon_kS_e\,\Psi_k(H),\label{eq:at-D}\\
    \bm C&=\int_{W_e}(\Rker*\bm\nu_e)(t)\,\dd t
      =S_e\int_0^H\bar\Rker(H-v)\,\bigl(\upsilon_k\psi_k(v)\bigr)_{k\in[M]}\,\dd v,\label{eq:at-C}\\
    \bm B&=\int_{W_e}\mathcal K(\bm\mu+\bm\nu_{-e})(t)\,\dd t=\bm B^{\mu}+\bm O,
    \qquad \bm B^{\mu}=\int_{W_e}\mathcal K\bm\mu,\quad \bm O=\int_{W_e}\mathcal K\bm\nu_{-e}.\label{eq:at-B}
  \end{align}
  In the cluster representation, $D_k$, $C_k$, $B^{\mu}_k$ and $O_k$ are the
  expected numbers of stream-$k$ incidents in $W_e$ of four kinds: immigrants of event
  $e$; descendants, of any generation, of those immigrants; members of clusters rooted
  in background immigrants; and members of clusters rooted in other events. The
  baseline $B_k$ is the expected count in $W_e$ under the same model with $S_e=0$, so
  $D_k+C_k$ is the expected excess attributable to event $e$.
  \item \emph{Totals.} As $H\to\infty$, $\bm D\uparrow\bm d_e$,
  $\bm D+\bm C\uparrow\Mmat\bm d_e$ and $\bm C\uparrow(\Mmat-\Imat)\bm d_e=\Amat\Mmat\bm d_e$.
  In the univariate case, $\Amat=n$ and $d_e>0$. The contagion share of the incidents
  attributable to the event then equals the branching ratio:
  \begin{equation}\label{eq:at-share-n}
    \frac{C(\infty)}{D(\infty)+C(\infty)}=n .
  \end{equation}
  \item \emph{Finite horizons.} For every $H$,
  $\bm C(H)\le\bar\Rker(H)\,\bm D(H)\le(\Mmat-\Imat)\,\bm D(H)$. In the univariate case,
  for every $H$ with $D(H)>0$,
  \[
  \frac{C(H)}{D(H)+C(H)}\le\frac{\bar\Rker(H)}{1+\bar\Rker(H)}\le n,
  \]
  and the ratio tends to $n$ as $H\to\infty$. If all streams share one response kernel,
  $\psi_k=\psi$, then for each stream $k$ the share
  $C_k(H)/\{D_k(H)+C_k(H)\}$ is at most its limit $C_k(\infty)/\{D_k(\infty)+C_k(\infty)\}$.
  \item \emph{Conditioning on the history before the event.} Let
  $\bm h_{\tau_e}(t)=\int_{[0,\tau_e]}\Phik(t-s)\,\dN(s)$, $t>\tau_e$, be the
  excitation carried over from incidents up to $\tau_e$. Then
  \begin{equation}\label{eq:at-cond}
    \E\bigl[N_k(W_e)\mid\Hist_{\tau_e}\bigr]=B^{\mid e}_k+D_k+C_k,\qquad
    \bm B^{\mid e}=\int_{W_e}\mathcal K_{\tau_e}\bigl(\bm\mu+\bm\nu_{-e}+\bm h_{\tau_e}\bigr)(t)\,\dd t ,
  \end{equation}
  with $D_k$ and $C_k$ as in (a). Conditioning on the past changes only the baseline.
\end{enumerate}
\end{theorem}

Part (b) gives the branching ratio a direct reading. In a univariate model, the
fraction of event-attributable incidents that are contagion is $n$, whatever the
excitation kernel, the response kernel $\psi$ and the size of the event: it is a property
of the cascade, not of the event. In a multivariate model the stream-$k$ share is
$1-d_{e,k}/(\Mmat\bm d_e)_k$, which depends on how the event loads on the other streams
and can exceed $\alpha_{kk}$ (\cref{sec:at-example}). Part (c) says that contagion takes
time, so shares reported at short horizons understate the eventual contagion. The scalar
ordering can fail when streams respond on different time scales. With
$\Amat=\bigl(\begin{smallmatrix}0.5&0\\0.4&0.2\end{smallmatrix}\bigr)$, all
$\beta_{kj}=3$, $\bm\eta=(5,0.2)$ and $\bm d_e=(10,5)^{\top}$, the share of stream 2 is
$0.879$ at $H=0.83$ but $0.692$ in the limit; the vector bound in (c) holds in every case.
This is one reason why the analysis plan reports several horizons
(\cref{sec:plan-estimands}). Part (d) assigns the excitation left by incidents before the
event to the baseline, so a surge that begins during an already elevated period is not
credited to the event. It is the version to use when decomposing the window after a
specific event, given the observed history.

Two supplements are in \cref{app:at-remarks}. A time-varying background is passed through
the filter $\delta_0\Imat+\Rker$, which amplifies slow variation by up to $\Mmat$; a calendar
effect that coincides with $\tau_e$ but is omitted from $\bm\mu$ is attributed to $\bm D$
and $\bm C$, a time-local analogue of \cref{thm:spurious} (\cref{rem:at-background}). And the
latent parent labels used by EM split a \emph{realised} surge into direct, contagion and
baseline parts whose expectations are the components of \eqref{eq:at-decomp}
(\cref{rem:at-expost}).

\subsection{Impulse responses for exponential kernels}\label{sec:at-irf}

For the exponential kernels of \cref{sec:th-hawkes}, the resolvent, the mean intensity
and the finite-horizon totals reduce to matrix exponentials.

\begin{proposition}[Impulse responses]\label{prop:irf}
\begin{enumerate}[(a)]
  \item \emph{Univariate.} If $M=1$ and $\varphi(t)=n\beta e^{-\beta t}$ with $n<1$,
  then
  \begin{equation}\label{eq:at-irf-uni}
    \Rker(t)=n\beta\,e^{-\beta(1-n)t},\qquad
    \bar\Rker(t)=\frac{n}{1-n}\bigl(1-e^{-\beta(1-n)t}\bigr).
  \end{equation}
  \item \emph{Decay rates set by the receiving stream.} Suppose
  $\varphi_{kj}(t)=\alpha_{kj}\beta_{k\bullet}e^{-\beta_{k\bullet}t}$, and let
  $\bm\Omega=\operatorname{diag}(\beta_{1\bullet},\dots,\beta_{M\bullet})$. Then
  \begin{equation}\label{eq:at-irf-row}
    \Rker(t)=e^{-\bm\Omega(\Imat-\Amat)t}\,\bm\Omega\Amat,
  \end{equation}
  and $-\bm\Omega(\Imat-\Amat)$ is Hurwitz, that is, all its eigenvalues have negative
  real part. The excitation part $\bm y=\Rker*\bm\nu$ of the mean intensity
  $\mathcal K\bm\nu=\bm\nu+\bm y$ solves the linear system
  $\bm y'=-\bm\Omega(\Imat-\Amat)\bm y+\bm\Omega\Amat\bm\nu$ with $\bm y(0)=\zero$.
  With a common rate, $\bm\Omega=\beta\Imat$, this gives
  $\Rker(t)=\beta\Amat e^{-\beta(\Imat-\Amat)t}$.
  \item \emph{General exponential kernels.} Suppose
  $\varphi_{kj}(t)=\alpha_{kj}\beta_{kj}e^{-\beta_{kj}t}$. Index $\mathbb R^{M^2}$ by
  pairs $(k,j)$, and define $\bm P:\mathbb R^M\to\mathbb R^{M^2}$,
  $\bm S:\mathbb R^{M^2}\to\mathbb R^M$ and $\bm Q:\mathbb R^{M^2}\to\mathbb R^{M^2}$ by
  \[
    (\bm P\bm v)_{kj}=\alpha_{kj}\beta_{kj}v_j,\qquad
    (\bm S\bm z)_k=\sum_{j}z_{kj},\qquad
    (\bm Q\bm z)_{kj}=-\beta_{kj}z_{kj}+\alpha_{kj}\beta_{kj}\sum_{l}z_{jl}.
  \]
  Then $\Rker(t)=\bm S\,e^{\bm Qt}\bm P$, and $\bm Q$ is Hurwitz. The coordinate
  $z_{kj}$ of $\bm z(t)=\int_0^te^{\bm Q(t-s)}\bm P\bm\nu(s)\,\dd s$ is the mean
  excitation of stream $k$ by stream-$j$ incidents.
  \item \emph{Finite-horizon totals for exponential shock responses.} Let
  $\psi_k(t)=\eta_ke^{-\eta_kt}$. For rates $\eta,\gamma>0$, let
  $\Pi_{\eta,\gamma}(H)=\Pr(X_\eta+X_\gamma\le H)$ be the hypoexponential distribution
  function. Here $X_\eta$ and $X_\gamma$ are independent exponential variables with
  rates $\eta$ and $\gamma$, and $\dot\Pi_{\eta,\gamma}$ denotes the density. Explicitly,
  \begin{equation}\label{eq:at-hypoexp}
    \Pi_{\eta,\gamma}(H)=
    \begin{cases}
      1-\dfrac{\gamma e^{-\eta H}-\eta e^{-\gamma H}}{\gamma-\eta}, & \eta\neq\gamma,\\[1.6ex]
      1-(1+\gamma H)\,e^{-\gamma H}, & \eta=\gamma\quad\text{(resonance)}.
    \end{cases}
  \end{equation}
  In the univariate case, let $d=\upsilon_1S_e$ and $\gamma=\beta(1-n)$. Then
  \begin{equation}\label{eq:at-DC-uni}
    D(H)=d\,\bigl(1-e^{-\eta H}\bigr),\qquad
    C(H)=\frac{n}{1-n}\,d\,\Pi_{\eta,\gamma}(H).
  \end{equation}
  The intensity attributable to the event at lag $t$ is
  $d\,\{\eta e^{-\eta t}+\tfrac{n}{1-n}\dot\Pi_{\eta,\gamma}(t)\}$. Its contagion
  part peaks at $t^*=\log(\eta/\gamma)/(\eta-\gamma)$, which equals $1/\gamma$ at
  resonance.

  If all kernels share one decay rate $\beta$ and all responses share one rate $\eta$, then
  \begin{equation}\label{eq:at-DC-common}
    \bm D(H)=\bigl(1-e^{-\eta H}\bigr)\bm d_e,\qquad
    \bm C(H)=(\Mmat-\Imat)\,\Pi_{\eta,\bm G}(H)\,\bm d_e,\qquad
    \bm G=\beta(\Imat-\Amat).
  \end{equation}
  Here
  $\Pi_{\eta,\bm G}(H)=\int_0^H\eta e^{-\eta v}\bigl(\Imat-e^{-\bm G(H-v)}\bigr)\,\dd v$
  is the matrix function induced by the entire map $x\mapsto\Pi_{\eta,x}(H)$.

  In general, with pair-specific $\beta_{kj}$ and stream-specific $\eta_k$, consider
  the state $\bm w=(\bm\nu_e,\bm z,\bm D,\bm C)$: the event's immigrant intensity, the
  excitation state of (c), and the two running totals. Then
  $(\bm D(H),\bm C(H))$ are the last two blocks of $\exp(H\bm Z)\,\bm w_0$, where
  \begin{equation}\label{eq:at-augmented}
    \bm Z=\begin{pmatrix}
      -\operatorname{diag}(\bm\eta) & \zero & \zero & \zero\\
      \bm P & \bm Q & \zero & \zero\\
      \Imat & \zero & \zero & \zero\\
      \zero & \bm S & \zero & \zero
    \end{pmatrix},
    \qquad
    \bm w_0=\begin{pmatrix}\operatorname{diag}(\bm\eta)\,\bm d_e\\ \zero\\ \zero\\ \zero\end{pmatrix}.
  \end{equation}
\end{enumerate}
\end{proposition}

Three features of \cref{prop:irf} matter for interpretation. First, the cascade decays
at rate $\gamma=\beta(1-n)$, more slowly than a single kernel, which matches the
correlation time in \cref{sec:th-second}. As $n\uparrow1$, $C(H)$ approaches its total
only slowly, so the horizon should exceed a few multiples of $\max\{1/\eta,1/\gamma\}$
before shares are compared across events. Second, in the univariate model the contagion
count up to $H$ is a fixed multiple of the distribution function of a sum of two delays:
from the event to a direct response (rate $\eta$), and from that response to a randomly
chosen descendant (rate $\gamma$, since $\Rker/\!\int\Rker$ is the $\mathrm{Exp}(\gamma)$
density). Third, at resonance, $\eta=\gamma$, the singularity in \eqref{eq:at-hypoexp} is
removable; the matrix exponential \eqref{eq:at-augmented} needs no special case
\citep{vanloan1978}, and the simulation code evaluates it directly. The background part
$\bm B^\mu$ is computed in the same way, with forcing $\bm\mu$ in place of $\bm\nu_e$.

\subsection{Online-to-offline spillover: a channel decomposition}\label{sec:at-channel}

Partition the streams into offline streams $\mathcal S_{\mathrm{off}}$ and online
streams $\mathcal S_{\mathrm{on}}$. Let $\bm E$ keep the online-to-offline entries of
$\Amat$, $E_{kj}=\alpha_{kj}$ for $k\in\mathcal S_{\mathrm{off}}$ and
$j\in\mathcal S_{\mathrm{on}}$, with all other entries zero. Set
$\Amat^\circ=\Amat-\bm E$ and $\Mmat^\circ=(\Imat-\Amat^\circ)^{-1}$. Write
$\bm C(H;\Amat)$ for the contagion component of \cref{thm:attribution} computed with
branching matrix $\Amat$, holding fixed every other parameter: the kernel shapes
$g_{kj}$, the background, the loadings and the response kernels. Only the
online-to-offline entries are set to zero. The online-to-online entries are kept;
zeroing the whole column of an online stream would also remove its excitation of
other online streams, and would answer a different question.

\begin{definition}[Online channel]\label{def:at-channel}
The \emph{online-channel contagion} of event $e$ over horizon $H$ is
\[
\bm C^{\mathrm{on}}(H)=\bm C(H;\Amat)-\bm C(H;\Amat^\circ).
\]
For an offline stream $k$ with $C_k(H)>0$, its \emph{online-channel share} is
$\sigma_k(H)=C^{\mathrm{on}}_k(H)/C_k(H)$.
\end{definition}

The direct component does not depend on $\Amat$. Hence $\bm C^{\mathrm{on}}$ is also the
reduction in all event-attributable incidents when the online-to-offline entries are
set to zero.

\begin{proposition}[Path sums and bounds]\label{prop:at-channel}
Under the standing assumptions:
\begin{enumerate}[(a)]
  \item \emph{Ancestral reading.} $C^{\mathrm{on}}_k(H)$ is the expected number of
  stream-$k$ incidents in $W_e$ that descend from an immigrant of event $e$ along an
  ancestral line that contains at least one online-to-offline transmission.
  \item \emph{Path sum.} $\spr(\Mmat^\circ\bm E)<1$ and
  \begin{equation}\label{eq:at-pathsum}
    \Mmat=\sum_{r\ge0}(\Mmat^\circ\bm E)^r\,\Mmat^\circ .
  \end{equation}
  Entry $(k,j)$ of the $r$th term is the total weight $\prod\alpha$ of the walks from
  $j$ to $k$ in the excitation graph with exactly $r$ online-to-offline steps.
  Consequently
  \begin{equation}\label{eq:at-chan-inf}
    \bm C^{\mathrm{on}}(\infty)=(\Mmat-\Mmat^\circ)\,\bm d_e
    =\Mmat\bm E\Mmat^\circ\bm d_e=\Mmat^\circ\bm E\Mmat\bm d_e .
  \end{equation}
  \item \emph{Bounds.} For every $H\in(0,\infty]$, $\zero\le\bm C^{\mathrm{on}}(H)\le\bm C(H)$,
  so $\sigma_k(H)\in[0,1]$. Moreover:
  \begin{itemize}
    \item $C^{\mathrm{on}}_k(\infty)>0$ if and only if some walk of positive weight
    runs from a stream $j$ with $d_{e,j}>0$ to $k$ and contains an online-to-offline
    step.
    \item $\bm C^{\mathrm{on}}(\infty)\ge\Mmat^\circ\bm E\Mmat^\circ\bm d_e$, which
    counts the walks with a single crossing.
    \item Equality holds when there is no offline-to-online excitation, that is,
    $\alpha_{kj}=0$ for all $k\in\mathcal S_{\mathrm{on}}$ and
    $j\in\mathcal S_{\mathrm{off}}$. Then, on the offline streams,
    \begin{equation}\label{eq:at-chan-nofb}
      \bm C^{\mathrm{on}}_{\mathrm{off}}(\infty)
      =(\Imat-\Amat_{\mathrm{off},\mathrm{off}})^{-1}\Amat_{\mathrm{off},\mathrm{on}}
      (\Imat-\Amat_{\mathrm{on},\mathrm{on}})^{-1}\bm d_{e,\mathrm{on}},
    \end{equation}
    which vanishes when the event has no direct online response.
  \end{itemize}
  \item \emph{Non-additivity.} Let $\mathcal E_1$ and $\mathcal E_2$ be disjoint sets of
  entries of $\Amat$, for instance the channels from two platforms. Let
  $\bm C^{\mathcal E}$ be defined as in \cref{def:at-channel}, with the entries in
  $\mathcal E$ set to zero. Then, for every $H$,
  \[
    \max\bigl(\bm C^{\mathcal E_1},\bm C^{\mathcal E_2}\bigr)\le\bm C^{\mathcal E_1\cup\mathcal E_2}
    \le\bm C^{\mathcal E_1}+\bm C^{\mathcal E_2}.
  \]
  The gap $\bm C^{\mathcal E_1}+\bm C^{\mathcal E_2}-\bm C^{\mathcal E_1\cup\mathcal E_2}$
  is the expected number of descendants whose ancestral line uses both channels.
\end{enumerate}
\end{proposition}

Part (d) means that channel shares are not additive. The shares of two platforms can sum
to more than the share of both together. Each share answers its own counterfactual
question. Shares should not be presented as slices of one pie.

\begin{remark}[What the channel decomposition does and does not show]\label{rem:at-causal}
$\bm C^{\mathrm{on}}$ is a \emph{model-based counterfactual}: the expected reduction in
event-attributable offline incidents if online-to-offline excitation were switched off and
nothing else changed. It becomes a causal statement, for instance about the effect of
moderation, only under further assumptions.
\begin{enumerate}[(i)]
  \item \emph{Excitation is influence.} By \cref{prop:granger},
  $\alpha_{\mathrm{off},\mathrm{on}}>0$ is equivalent to Granger causality, which is
  predictive, not causal. Cross-excitation sometimes leaves a second-order trace that no
  common-shock model reproduces (\cref{ex:th-mv-cross}), but an omitted common driver can
  still create apparent excitation (\cref{thm:spurious}, \cref{rem:th-granger-caveats}(i)).
  \item \emph{Modularity.} An intervention must remove exactly the entries in $\bm E$ and
  leave everything else unchanged; interventions on platforms may displace activity rather
  than remove it.
  \item \emph{Measurement and timestamps.} The online stream is a classifier output and
  offline streams are thinned, so the decomposition should be computed on the true scale
  (\cref{rem:at-thinning}); and report-time stamps can make online reactions precede the
  record of an incident (\cref{rem:th-granger-caveats}(iv)).
\end{enumerate}
In line with the analysis plan (\cref{sec:plan-decisions}), $\bm C^{\mathrm{on}}$ and $\sigma_k$
are reported as model-based quantities and are not described as causal effects.
\end{remark}

\subsection{Uncertainty}\label{sec:at-uncertainty}

Every component above is a smooth function $g(\vartheta)$ of the parameter vector
$\vartheta$ on the region $\spr(\Amat)<1$. Given posterior or parametric-bootstrap draws
$\vartheta^{(1)},\dots,\vartheta^{(R)}$, we compute $(\bm B,\bm D,\bm C,\bm C^{\mathrm{on}})^{(r)}$
exactly for each draw and report quantiles. This propagates the dependence between
components and respects the constraints $\sigma_k\in[0,1]$ and, in univariate models,
$C/(D+C)\le n$; infinite-horizon totals are reported only for draws with
$\spr(\Amat^{(r)})<1$. The analysis plan uses this route. The delta method, with the
closed-form gradients of \cref{app:at-grad}, is a cheaper alternative, but it is unreliable
near the boundary $\alpha_{kj}=0$ (\cref{sec:est-lrt}) and near criticality, where $\Mmat$
has a pole.

\begin{remark}[Under-recording]\label{rem:at-thinning}
Suppose each stream-$k$ incident is recorded independently with probability $p_k$. Then
the expected number of recorded incidents in $W_e$ is $p_k(B_k+D_k+C_k)$. Recording
scales all three components of a stream by the same factor. The within-stream shares
$C_k/(D_k+C_k)$ and $\sigma_k$ are therefore invariant to $p_k$, provided they are
computed from true-scale parameters.

Fitting the recorded streams as if they were complete does not have this invariance. By
\cref{thm:thinning}, moment-based estimators of the branching ratio from a recorded
univariate stream converge to $h_p(n)<n$. The contagion share implied by such a fit,
which equals the fitted branching ratio by \eqref{eq:at-share-n}, understates
contagion. The dark-figure correction (\cref{cor:th-darkfig}) should therefore precede
attribution.
\end{remark}

\Cref{sec:at-example} works through a two-stream example with invented parameters, in
which the offline contagion share is about twice $\alpha_{\mathrm{off},\mathrm{off}}$ and
most of the offline contagion passes through the online stream.
